\documentclass[10pt,a4paper]{amsart}
\usepackage{amsmath,amssymb,amsthm,graphicx}
\usepackage{verbatim}
\newenvironment{reference}{\comment}{\endcomment}
%\newenvironment{reference}{}{}
\newenvironment{slogan}{\comment}{\endcomment}
\newenvironment{history}{\comment}{\endcomment}

% For commutative diagrams we use Xy-pic
\usepackage[all]{xy}

% We use 2cell for 2-commutative diagrams.
\xyoption{2cell}
% Theorem environments.
%
\theoremstyle{plain}
\newtheorem{theorem}[subsection]{Theorem}
\newtheorem{proposition}[subsection]{Proposition}
\newtheorem{lemma}[subsection]{Lemma}

\theoremstyle{definition}
\newtheorem{definition}[subsection]{Definition}
\newtheorem{example}[subsection]{Example}
\newtheorem{exercise}[subsection]{Exercise}
\newtheorem{situation}[subsection]{Situation}

\theoremstyle{remark}
\newtheorem{remark}[subsection]{Remark}
\newtheorem{remarks}[subsection]{Remarks}

\numberwithin{equation}{subsection}

% Macros
%
\def\lim{\mathop{\mathrm{lim}}\nolimits}
\def\colim{\mathop{\mathrm{colim}}\nolimits}
\def\Spec{\mathop{\mathrm{Spec}}}
\def\Hom{\mathop{\mathrm{Hom}}\nolimits}
\def\Ext{\mathop{\mathrm{Ext}}\nolimits}
\def\SheafHom{\mathop{\mathcal{H}\!\mathit{om}}\nolimits}
\def\SheafExt{\mathop{\mathcal{E}\!\mathit{xt}}\nolimits}
\def\Sch{\mathit{Sch}}
\def\Mor{\mathop{\mathrm{Mor}}\nolimits}
\def\Ob{\mathop{\mathrm{Ob}}\nolimits}
\def\Sh{\mathop{\mathit{Sh}}\nolimits}
\def\NL{\mathop{N\!L}\nolimits}
\def\CH{\mathop{\mathrm{CH}}\nolimits}
\def\proetale{{pro\text{-}\acute{e}tale}}
\def\etale{{\acute{e}tale}}
\def\QCoh{\mathit{QCoh}}
\def\Ker{\mathop{\mathrm{Ker}}}
\def\Im{\mathop{\mathrm{Im}}}
\def\Coker{\mathop{\mathrm{Coker}}}
\def\Coim{\mathop{\mathrm{Coim}}}

% Boxtimes
%
\DeclareMathSymbol{\boxtimes}{\mathbin}{AMSa}{"02}

%
% Macros for moduli stacks/spaces
%
\def\QCohstack{\mathcal{QC}\!\mathit{oh}}
\def\Cohstack{\mathcal{C}\!\mathit{oh}}
\def\Spacesstack{\mathcal{S}\!\mathit{paces}}
\def\Quotfunctor{\mathrm{Quot}}
\def\Hilbfunctor{\mathrm{Hilb}}
\def\Curvesstack{\mathcal{C}\!\mathit{urves}}
\def\Polarizedstack{\mathcal{P}\!\mathit{olarized}}
\def\Complexesstack{\mathcal{C}\!\mathit{omplexes}}
% \Pic is the operator that assigns to X its picard group, usage \Pic(X)
% \Picardstack_{X/B} denotes the Picard stack of X over B
% \Picardfunctor_{X/B} denotes the Picard functor of X over B
\def\Pic{\mathop{\mathrm{Pic}}\nolimits}
\def\Picardstack{\mathcal{P}\!\mathit{ic}}
\def\Picardfunctor{\mathrm{Pic}}
\def\Deformationcategory{\mathcal{D}\!\mathit{ef}}
\usepackage{cleveref}
\crefname{originallemma}{original lemma}{original lemmas}
\Crefname{originallemma}{Original lemma}{Original lemmas}
\makeatletter
\newlabel{lemma-localized-chern-pre}{{0F9K}{source}{}{sourcepointer-1}{}}
\newlabel{lemma-localized-chern-pre@cref}{{[originallemma][0][]0F9K}{[1][0][]source}{}{}{}}
\makeatother
\crefname{originalreference}{original reference}{original references}
\Crefname{originalreference}{Original reference}{Original references}
\makeatletter
\newlabel{remark-ring-loc-classes}{{0FA0}{source}{}{sourcepointer-2}{}}
\newlabel{remark-ring-loc-classes@cref}{{[originalreference][0][]0FA0}{[1][0][]source}{}{}{}}
\makeatother
\crefname{originallemma}{original lemma}{original lemmas}
\Crefname{originallemma}{Original lemma}{Original lemmas}
\makeatletter
\newlabel{lemma-additivity-on-perfect}{{0F9F}{source}{}{sourcepointer-3}{}}
\newlabel{lemma-additivity-on-perfect@cref}{{[originallemma][0][]0F9F}{[1][0][]source}{}{}{}}
\makeatother
\crefname{originallemma}{original lemma}{original lemmas}
\Crefname{originallemma}{Original lemma}{Original lemmas}
\makeatletter
\newlabel{lemma-silly-sum-c}{{0FAL}{source}{}{sourcepointer-4}{}}
\newlabel{lemma-silly-sum-c@cref}{{[originallemma][0][]0FAL}{[1][0][]source}{}{}{}}
\makeatother
\crefname{originalreference}{original reference}{original references}
\Crefname{originalreference}{Original reference}{Original references}
\makeatletter
\newlabel{remark-res-push}{{0FA1}{source}{}{sourcepointer-5}{}}
\newlabel{remark-res-push@cref}{{[originalreference][0][]0FA1}{[1][0][]source}{}{}{}}
\makeatother
\crefname{originallemma}{original lemma}{original lemmas}
\Crefname{originallemma}{Original lemma}{Original lemmas}
\makeatletter
\newlabel{lemma-silly}{{0F9H}{source}{}{sourcepointer-6}{}}
\newlabel{lemma-silly@cref}{{[originallemma][0][]0F9H}{[1][0][]source}{}{}{}}
\makeatother
\crefname{originallemma}{original lemma}{original lemmas}
\Crefname{originallemma}{Original lemma}{Original lemmas}
\makeatletter
\newlabel{lemma-homomorphism}{{0FAV}{source}{}{sourcepointer-7}{}}
\newlabel{lemma-homomorphism@cref}{{[originallemma][0][]0FAV}{[1][0][]source}{}{}{}}
\makeatother
\crefname{originallemma}{original lemma}{original lemmas}
\Crefname{originallemma}{Original lemma}{Original lemmas}
\makeatletter
\newlabel{lemma-cap-commutative-chern}{{02UA}{source}{}{sourcepointer-8}{}}
\newlabel{lemma-cap-commutative-chern@cref}{{[originallemma][0][]02UA}{[1][0][]source}{}{}{}}
\makeatother
\title{The total localized Chern class of a direct sum is the product of the total classes of the summands}
\author{The Stacks Project Authors}
\date{stacks-project:chow.tex:local-snapshot-20261012; GNU FDL 1.2 or later}
\begin{document}
\maketitle
\noindent\textit{Editorial note.} All delivered mathematical text is the author's own, copied from the pinned Stacks Project chapter file chow.tex; no statement or proof step was rewritten, reordered or removed.\par\medskip
\noindent\textit{Editorial note.} Preamble provenance: the pinned chapter file opens with an input of the shared preamble and carries no preamble text of its own, so the project's single Stacks-wide preamble staged at sources/stacks-chow/preamble.tex is excerpted verbatim as two spans, lines 14-18 (Xy-pic setup) and lines 154-217 (the native theorem declarations and mathematical macros). Nothing was added, removed or redefined; the wrapper class and the amsmath/amssymb/amsthm/graphicx base packages are supplied by the assembler.\par\medskip
\noindent\textit{Editorial note.} The counted statement and proof contain no cross reference, no citation and no included file, so the delivered body is self contained.\par\medskip
\noindent\textit{Editorial note.} Canonical location of the counted proposition: Stacks tag 0FAY (sources/stacks-chow/pinned-original-source.tex line 10499).\par\medskip
\noindent\textit{Allowed standard prerequisites (editorial).}\begin{itemize}
\item The localized Chern classes and the truncated total class c\textasciicircum{}\{(p)\}(Q). Same chapter, Lemma labelled lemma-localized-chern-pre, cited at proof lines 10519-10521 and 10536-10540, and Remark labelled remark-ring-loc-classes, invoked in the statement. Fixes the classes being multiplied and the reduction of c'\_i(Q) to c\_i of the restriction to the fiber over zero.
\item Additivity of the localized Chern classes. Same chapter, Lemma labelled lemma-additivity-on-perfect, cited at proof lines 10525-10527. Handles the degrees below p + p'.
\item The product formula over E and the description of c'\_p via pushforward and C. Same chapter, Lemmas labelled lemma-silly-sum-c and lemma-silly, cited at proof lines 10532-10534 and 10565-10567, and Lemma labelled lemma-localized-chern-pre for the formula involving C. Lets the computation be carried out on the closed subscheme E inside W\_infinity.
\item Commutation of bivariant Chern classes with capping and the homomorphism lemma. Same chapter, Lemma labelled lemma-cap-commutative-chern, cited at proof lines 10592-10594, and Lemma labelled lemma-homomorphism, cited at proof lines 10575-10578 and 10586-10589. Settles the two mixed cases where one index is below the truncation bound.
\end{itemize}\medskip
\ifdefined\hypertarget\hypertarget{sourcepointer-1}{}\fi
\noindent\textit{Source reference.} Original lemma 0F9K: sources/stacks-chow/pinned-original-source.tex, source lines 10204--10204. Same-chapter pointer into the author environment lemma at source line 10204; invoked at proof line 10522 of the counted proof. It is retained only as a navigation pointer; every load-bearing step is inside the delivered span.\par\medskip
\ifdefined\hypertarget\hypertarget{sourcepointer-2}{}\fi
\noindent\textit{Source reference.} Original reference 0FA0: sources/stacks-chow/pinned-original-source.tex, source lines 6419--6419. Same-chapter pointer into the author environment remark at source line 6419; invoked at proof line 10547 of the counted proof. It is retained only as a navigation pointer; every load-bearing step is inside the delivered span.\par\medskip
\ifdefined\hypertarget\hypertarget{sourcepointer-3}{}\fi
\noindent\textit{Source reference.} Original lemma 0F9F: sources/stacks-chow/pinned-original-source.tex, source lines 9427--9427. Same-chapter pointer into the author environment lemma at source line 9427; invoked at proof line 10524 of the counted proof. It is retained only as a navigation pointer; every load-bearing step is inside the delivered span.\par\medskip
\ifdefined\hypertarget\hypertarget{sourcepointer-4}{}\fi
\noindent\textit{Source reference.} Original lemma 0FAL: sources/stacks-chow/pinned-original-source.tex, source lines 9889--9889. Same-chapter pointer into the author environment lemma at source line 9889; invoked at proof line 10535 of the counted proof. It is retained only as a navigation pointer; every load-bearing step is inside the delivered span.\par\medskip
\ifdefined\hypertarget\hypertarget{sourcepointer-5}{}\fi
\noindent\textit{Source reference.} Original reference 0FA1: sources/stacks-chow/pinned-original-source.tex, source lines 6451--6451. Same-chapter pointer into the author environment remark at source line 6451; invoked at proof line 10554 of the counted proof. It is retained only as a navigation pointer; every load-bearing step is inside the delivered span.\par\medskip
\ifdefined\hypertarget\hypertarget{sourcepointer-6}{}\fi
\noindent\textit{Source reference.} Original lemma 0F9H: sources/stacks-chow/pinned-original-source.tex, source lines 9695--9695. Same-chapter pointer into the author environment lemma at source line 9695; invoked at proof line 10561 of the counted proof. It is retained only as a navigation pointer; every load-bearing step is inside the delivered span.\par\medskip
\ifdefined\hypertarget\hypertarget{sourcepointer-7}{}\fi
\noindent\textit{Source reference.} Original lemma 0FAV: sources/stacks-chow/pinned-original-source.tex, source lines 10353--10353. Same-chapter pointer into the author environment lemma at source line 10353; invoked at proof line 10571 of the counted proof. It is retained only as a navigation pointer; every load-bearing step is inside the delivered span.\par\medskip
\ifdefined\hypertarget\hypertarget{sourcepointer-8}{}\fi
\noindent\textit{Source reference.} Original lemma 02UA: sources/stacks-chow/pinned-original-source.tex, source lines 7581--7581. Same-chapter pointer into the author environment lemma at source line 7581; invoked at proof line 10594 of the counted proof. It is retained only as a navigation pointer; every load-bearing step is inside the delivered span.\par\medskip
\begin{lemma}
\label{lemma-localized-chern-pre-sum-c}
In Lemma \ref{lemma-localized-chern-pre} assume $Q|_T$ is isomorphic
to a finite locally free $\mathcal{O}_T$-module of rank $< p$.
Assume we have another perfect object $Q' \in D(\mathcal{O}_W)$
whose Chern classes are defined with $Q'|_T$ isomorphic to a
finite locally free $\mathcal{O}_T$-module of rank $< p'$ placed
in cohomological degree $0$. With notation as in
Remark \ref{remark-ring-loc-classes} set
$$
c^{(p)}(Q) = 1 + c_1(Q|_{X \times \{0\}}) + \ldots +
c_{p - 1}(Q|_{X \times \{0\}}) +
c'_{p}(Q) + c'_{p + 1}(Q) + \ldots
$$
in $A^{(p)}(Z \to X)$ with $c'_i(Q)$ for $i \geq p$ as in
Lemma \ref{lemma-localized-chern-pre}. Similarly for $c^{(p')}(Q')$ and
$c^{(p + p')}(Q \oplus Q')$.
Then $c^{(p + p')}(Q \oplus Q') = c^{(p)}(Q)c^{(p')}(Q')$
in $A^{(p + p')}(Z \to X)$.
\end{lemma}

\begin{proof}
Recall that the image of $c'_i(Q)$ in $A^p(X)$ is equal to
$c_i(Q|_{X \times \{0\}})$ for $i \geq p$ and similarly for
$Q'$ and $Q \oplus Q'$, see Lemma \ref{lemma-localized-chern-pre}.
Hence the equality in degrees $< p + p'$ follows from the
additivity of Lemma \ref{lemma-additivity-on-perfect}.

\medskip\noindent
Let's take $n \geq p + p'$.
As in the proof of Lemma \ref{lemma-localized-chern-pre}
let $E \subset W_\infty$ denote the inverse image of $Z$.
Observe that we have the equality
$$
c^{(p + p')}(Q|_E \oplus Q'|_E) =
c^{(p)}(Q|_E)c^{(p')}(Q'|_E)
$$
in $A^{(p + p')}(E \to W_\infty)$ by Lemma \ref{lemma-silly-sum-c}.
Since by construction
$$
c'_p(Q \oplus Q') = (E \to Z)_* \circ c'_p(Q|_E \oplus Q'|_E) \circ C
$$
we conclude that suffices to show for all $i + j = n$ we have
$$
(E \to Z)_* \circ c^{(p)}_i(Q|_E)c^{(p')}_j(Q'|_E) \circ C
=
c^{(p)}_i(Q)c^{(p')}_j(Q')
$$
in $A^n(Z \to X)$ where the multiplication is the one from
Remark \ref{remark-ring-loc-classes} on both sides. There are
three cases, depending on whether $i \geq p$, $j \geq p'$, or both.

\medskip\noindent
Assume $i \geq p$ and $j \geq p'$. In this case the products are
defined by inserting $(E \to W_\infty)_*$, resp.\ $(Z \to X)_*$ in between
the two factors and taking compositions as bivariant classes, see
Remark \ref{remark-res-push}.
In other words, we have to show
$$
(E \to Z)_* \circ c'_i(Q|_E) \circ
(E \to W_\infty)_* \circ c'_j(Q'|_E) \circ C =
c'_i(Q) \circ (Z \to X)_* \circ c'_j(Q')
$$
By Lemma \ref{lemma-silly} the left hand side is equal to
$$
(E \to Z)_* \circ c'_i(Q|_E) \circ c_j(Q'|_{W_\infty}) \circ C
$$
Since $c'_i(Q) = (E \to Z)_* \circ c'_i(Q|_E) \circ C$
the right hand side is equal to
$$
(E \to Z)_* \circ c'_i(Q|_E) \circ C \circ (Z \to X)_* \circ c'_j(Q')
$$
which is immediately seen to be equal to the above
by Lemma \ref{lemma-homomorphism}.

\medskip\noindent
Assume $i \geq p$ and $j < p$. Unwinding the products
in this case we have to show
$$
(E \to Z)_* \circ c'_i(Q|_E) \circ c_j(Q'|_{W_\infty}) \circ C =
c'_i(Q) \circ c_j(Q'|_{X \times \{0\}})
$$
Again using that $c'_i(Q) = (E \to Z)_* \circ c'_i(Q|_E) \circ C$
we see that it suffices to show $c_j(Q'|_{W_\infty}) \circ C =
C \circ c_j(Q'|_{X \times \{0\}})$ which is part of
Lemma \ref{lemma-homomorphism}.

\medskip\noindent
Assume $i < p$ and $j \geq p'$. Unwinding the products
in this case we have to show
$$
(E \to Z)_* \circ c_i(Q|_E) \circ c'_j(Q'|_E) \circ C =
c_i(Q|_{Z \times \{0\}}) \circ c'_j(Q')
$$
However, since $c'_j(Q|_E)$ and $c'_j(Q')$ are
bivariant classes, they commute with capping with Chern classes
(Lemma \ref{lemma-cap-commutative-chern}). Hence it suffices to prove
$$
(E \to Z)_* \circ c'_j(Q'|_E) \circ c_i(Q|_{W_\infty}) \circ C =
c'_j(Q') \circ c_i(Q|_{X \times \{0\}})
$$
which we reduces us to the case discussed in the preceding paragraph.
\end{proof}

\end{document}
